\documentclass{beamer}
\usetheme{Madrid}
\usecolortheme{default}
\usepackage{listings}
\usepackage{xcolor}
\usepackage{ctex}
\setmonofont{DejaVu Sans Mono}

% Code listing settings
\lstset{
    language=C++,
    basicstyle=\footnotesize\ttfamily,
    keywordstyle=\color{blue},
    commentstyle=\color{green!60!black},
    stringstyle=\color{red},
    numbers=left,
    numberstyle=\footnotesize\color{gray},
    stepnumber=1,
    numbersep=5pt,
    backgroundcolor=\color{white},
    showspaces=false,
    showstringspaces=false,
    showtabs=false,
    frame=single,
    tabsize=2,
    breaklines=true,
    breakatwhitespace=false
}

\title{Lab3 Sharing}
\author{李鹏达}
\date{\today}

\begin{document}

\frame{\titlepage}

% Slide 1: weightedRand
\begin{frame}[fragile]
\frametitle{Weighted Random Selection}
\begin{lstlisting}
int weightedRand(const std::vector<double> &Weights) {
  if (Weights.empty()) return 0;
  
  static std::mt19937 gen(static_cast<unsigned>(rand()));
  std::discrete_distribution<int> dist(Weights.begin(), 
                                        Weights.end());
  return dist(gen);
}
\end{lstlisting}
\end{frame}

% Slide 2: selectInput
\begin{frame}[fragile]
\frametitle{Input Selection Strategy}
\begin{lstlisting}
std::string selectInput(RunInfo Info) {
  if (SeedInputs.empty()) 
    return "";
  int Index = weightedRand(SeedWeight);
  MutationState = Index;
  return SeedInputs[Index];
}
\end{lstlisting}
\end{frame}

% Slide 3: selectMutationFn
\begin{frame}[fragile]
\frametitle{Mutation Function Selection}
\begin{lstlisting}
MutationFn *selectMutationFn(RunInfo &Info) {
  int Strat = weightedRand(MutationFnWeights);
  StrategyState = Strat;
  return MutationFns[Strat];
}
\end{lstlisting}
\end{frame}

% Slide 4: Mutation 0 and 1
\begin{frame}[fragile]
\frametitle{Mutation Strategies (1/3)}
\begin{lstlisting}
// Mutation 0: No mutation
std::string mutation_0(std::string Original) { 
  return Original; 
}

// Mutation 1: Insert random character
std::string mutation_1(std::string Original) {
  if (Original.length() <= 0)
    return Original;
  int Index = rand() % Original.length();
  return Original.insert(Index, 1, 
                         ALPHA[rand() % LENGTH_ALPHA]);
}
\end{lstlisting}
\end{frame}

% Slide 5: Mutation 2 and 3
\begin{frame}[fragile]
\frametitle{Mutation Strategies (2/3)}
\begin{lstlisting}
// Mutation 2: Swap adjacent characters
std::string mutation_2(std::string Original) {
  if (Original.length() <= 1) return Original;
  int Index = rand() % Original.length();
  int Addition = rand() % 2 == 0 ? 1 : -1;
  int SwapIndex = Index + Addition;
  if (SwapIndex < 0 || 
      static_cast<std::size_t>(SwapIndex) >= Original.length()) {
    SwapIndex = Index - Addition;
  }
  if (SwapIndex < 0 || 
      static_cast<std::size_t>(SwapIndex) >= Original.length()) {
    return Original;
  }
  char Temp = Original[Index];
  Original[Index] = Original[SwapIndex];
  Original[SwapIndex] = Temp;
  return Original;
}
\end{lstlisting}
\end{frame}

% Slide 6: Mutation 3 continued
\begin{frame}[fragile]
\frametitle{Mutation Strategies (2/3 continued)}
\begin{lstlisting}
// Mutation 3: Increment/decrement character
std::string mutation_3(std::string Original) {
  if (Original.length() <= 0)
    return Original;
  int Index = rand() % Original.length();
  int Addition = rand() % 2 == 0 ? 1 : -1;
  int AlphaIndex = 0;
  for (int i = 0; i < LENGTH_ALPHA; ++i) {
    if (ALPHA[i] == Original[Index]) {
      AlphaIndex = i;
      break;
    }
  }
  AlphaIndex = (AlphaIndex + Addition + LENGTH_ALPHA) 
               % LENGTH_ALPHA;
  Original[Index] = ALPHA[AlphaIndex];
  return Original;
}
\end{lstlisting}
\end{frame}

% Slide 7: Mutation 4 and 5
\begin{frame}[fragile]
\frametitle{Mutation Strategies (3/3)}
\begin{lstlisting}
// Mutation 4: Delete character
std::string mutation_4(std::string Original) {
  if (Original.length() <= 0)
    return Original;
  int Index = rand() % Original.length();
  return Original.erase(Index, 1);
}

// Mutation 5: Insert newline character
std::string mutation_5(std::string Original) {
  int Pos = Original.empty() ? 0 : 
            (rand() % (Original.size() + 1));
  char Special = '\n';
  Original.insert(Pos, 1, Special);
  return Original;
}
\end{lstlisting}
\end{frame}

% Slide 8: Feedback (Part 1)
\begin{frame}[fragile]
\frametitle{Feedback Algorithm (1/2)}
\begin{lstlisting}
void feedBack(std::string &Target, RunInfo &Info) {
  std::vector<std::string> RawCoverageData;
  readCoverageFile(Target, RawCoverageData);

  PrevCoverageState = CoverageState;
  CoverageState.clear();
  CoverageState.assign(RawCoverageData.begin(),
                       RawCoverageData.end());

  int Diff = coverageDiff(PrevCoverageState, CoverageState);

  if (Diff > 0) {
    SeedInputs.push_back(Info.MutatedInput);
    SeedWeight.push_back(10.0);
    
    SeedWeight[MutationState] *= 1.5;
    if (SeedWeight[MutationState] > 100.0) {
      SeedWeight[MutationState] = 100.0;
    }
\end{lstlisting}
\end{frame}

% Slide 9: Feedback (Part 2)
\begin{frame}[fragile]
\frametitle{Feedback Algorithm (2/2)}
\begin{lstlisting}
    MutationFnWeights[StrategyState] *= 1.5;
    if (MutationFnWeights[StrategyState] > 100.0) {
      MutationFnWeights[StrategyState] = 100.0;
    }
  } else if (Info.Passed) {
    if (MutationState >= 0 && 
        MutationState < static_cast<int>(SeedWeight.size())) {
      SeedWeight[MutationState] *= 0.9;
      if (SeedWeight[MutationState] < 0.1) {
        SeedWeight[MutationState] = 0.1;
      }
    }
    
    MutationFnWeights[StrategyState] *= 0.9;
    if (MutationFnWeights[StrategyState] < 0.1) {
      MutationFnWeights[StrategyState] = 0.1;
    }
  }
}
\end{lstlisting}
\end{frame}

\begin{frame}[fragile]
\frametitle{Why Mutation 5?}
Use mutation\_1:
$$
\frac{1}{6} \times \frac{1}{26} \times \frac{1}{2} =\frac{1}{312}
$$
100 times to get a newline that meets the condition:
$$
\left(1 - \left(1 - \frac{1}{312}\right)^{100}\right) \approx 0.27
$$
\end{frame}

\begin{frame}
\begin{center}
    \huge THANKS
\end{center}
\end{frame}

\end{document}